\section{Applications}

\input{sections/ai_scientists}

\subsection{Existing GPMs for Chemical Science}

The development of \glspl{gpm} for chemical science represents a departure from traditional single-task approaches. Rather than fine-tuning pre-trained models for specific tasks such as property prediction or molecular generation, these chemistry-aware models are intentionally designed to be capable of performing different chemical tasks. This multitask learning unifies related chemical tasks, boosting data efficiency and enabling emergent capabilities through joint training across domains. \Cref{tab:existing_gpms} shows an overview of the existing \glspl{gpm} in the chemical sciences.

\begin{table}[htbp]
\centering
\setlength{\tabcolsep}{5pt}
\renewcommand{\arraystretch}{1.3}
\setlength{\extrarowheight}{0.3ex}
\caption{\textbf{Non-comprehensive list of existing \glspl{gpm} in chemical sciences.} Most current models are trained on a mixture of general and domain-specific corpora. Models are sorted by the order they appear in the main text.}
\label{tab:existing_gpms}
\footnotesize
\begin{tabularx}{\linewidth}{
  l
  >{\raggedright\arraybackslash}X
  >{\raggedright\arraybackslash}X
  >{\raggedright\arraybackslash}X
  >{\raggedright\arraybackslash}X
}
\toprule
\textbf{Model} & \textbf{Architecture} & \textbf{Training} & \textbf{Representative Tasks} & \textbf{Representative Systems} \\
\midrule
\modelname{DARWIN 1.5} \autocite{xie2025darwin} & decoder-only transformer (LLaMA) with 7B parameters & fine-tuning on \glspl{qa}; multi-task fine-tuning on language-interfaced classification and regression tasks &  materials property prediction from natural language prompts & inorganic materials \\
\addlinespace[2.5pt]
\modelname{ChemLLM} \autocite{zhang2024chemllm} & decoder-only transformer (InternLM2) with 7B parameters & instruction-tuning on template-based \glsps{qa} using \gls{sft} or \gls{dpo} & multi-topic chemistry \gls{qa} & organic molecules \& reactions \\
\addlinespace[2.5pt]
\modelname{nach0} \autocite{livne2024nach0} &
encoder–decoder transformer (T5) with 250M or 780M parameters &
pre-trained using \gls{ssl} on scientific literature and \gls{smiles} strings; instruction-tuned &
natural language to \gls{smiles}; small-molecule property prediction &
organic molecules \& reactions \\
\addlinespace[2.5pt]
% \modelname{MatBERT} \autocite{qu2023leveraging, wan2024tokens} &
% encoder-only transformer (BERT) with 110M parameters &
% masked language modeling pre-trained on materials literature; task-specific fine-tuning &
% create embeddings for property prediction &
% inorganic materials \\
% \addlinespace[2.5pt]
\modelname{LLaMat} \autocite{mishra2024foundational} &
decoder-only transformer (LLaMA) with 7B parameters &
continued pre-training on materials literature and crystallographic data &
materials information extraction; understanding and generating \glspl{cif} &
inorganic materials \\
\addlinespace[2.5pt]
\modelname{ChemDFM} \autocite{zhao2024chemdfm} &
decoder-only transformer (LLaMA) with 13B parameters &
pre-trained on 34B tokens from chemistry papers and textbooks; instruction-tuned with \gls{qa} and \gls{smiles} &
\gls{smiles} to text; molecule/reaction property prediction &
organic molecules \& reactions \\
\addlinespace[2.5pt]
\modelname{ether0} \autocite{narayanan2025training} &
decoder-only transformer (Mistral-Small) with 24B parameters; \gls{moe} &
\gls{sft} on reasoning traces, then \gls{rl} &
molecular editing; one-step retrosynthesis; reaction prediction &
organic molecules \& reactions \\
\bottomrule
\end{tabularx}
\end{table}

\paragraph{Domain Pre-Training and Multitask Learning} \modelname{DARWIN 1.5} used a multitask approach to fine-tune \modelname{Llama-7B} through a two-stage process \autocite{xie2025darwin}. In the first step, the base model was fine-tuned on $332k$ scientific \gls{qa} pairs to establish foundational scientific reasoning. Subsequently, the model underwent multitask learning on $22$ different regression and classification tasks based on experimental datasets. \modelname{DARWIN 1.5}'s core idea is the use of \gls{lift} (see \Cref{sec:model_adaptation}) on diverse materials tasks to induce cross-task transfer during training. However, in some cases, the results revealed negative task-interaction, i.e., some tasks had diminished performance under multi-tasking.

\modelname{ChemLLM} followed a similar approach to \modelname{DARWIN 1.5}: template-based instruction tuning (ChemData) on $~7M$ \gls{qa} pairs. \autocite{zhang2024chemllm}.

While the above examples focus on combining different in-domain tasks, \modelname{nach0} coupled natural language with chemical data \autocite{livne2024nach0}, based on a unified encoder-decoder transformer architecture. The model was pre-trained using \gls{ssl} on a combination of \gls{smiles} strings and natural language from scientific literature, and then instruction-tuned on chemistry and \gls{nlp} tasks. This allows \modelname{nach0} to translate between natural language and \gls{smiles}, in addition to tasks like \gls{qa}, information extraction, and molecule/reaction generation.

% \textcite{wan2024tokens} developed a language-driven materials discovery framework that uses embeddings from \modelname{MatBERT}\autocite{qu2023leveraging}---a \gls{mlm} pre-trained on materials-science literature---to represent and generate novel crystal structures. Candidates were first recalled via similarity in the embedding space, then ranked using a multitask multi-gate \gls{moe} model that predicts the desired properties jointly. 
% Their method successfully identified novel high-performance materials (e.g., halide perovskites), and demonstrated that language representations encode latent knowledge for task-agnostic materials design.
\modelname{LLaMat} employed \gls{peft} to continue the pre-training on crystal structure data in \gls{cif} format, enabling the generation of thermodynamically stable structures \autocite{mishra2024foundational}.

\modelname{ChemDFM} scaled this concept significantly, implementing domain pre-training on over 34 billion tokens from chemical textbooks and research articles \autocite{zhao2024chemdfm}. After that, through comprehensive instruction tuning, \modelname{ChemDFM} familiarizes itself with chemical notation and patterns, distinguishing it from more materials-focused approaches like \modelname{LLaMat} through its broader chemical knowledge base.

\paragraph{Reasoning-Based Approaches} A recent development in chemical \glspl{gpm} incorporates explicit reasoning capabilities. The \modelname{ether0} model demonstrated this approach as a 24 billion-parameter reasoning model trained on over 640k chemically-verifiable problems (e.g., through code) across 375 tasks, including single-step retrosynthesis, molecular editing, and reaction prediction \autocite{narayanan2025training}. 
Unlike previous models, \modelname{ether0}'s training used a novel \gls{rl} approach (see \Cref{sec:rl}). First, \modelname{DeepSeek-R1} was prompted to create long \gls{cot} traces ending in a \gls{smiles}. These traces were used to fine-tune a pre-trained \modelname{Mistral-Small-24B} using \gls{sft}, resulting in a \enquote{base reasoner} model. This model underwent \gls{rl} using \gls{grpo} on several verifiable chemical tasks to create \enquote{specialist} models---one for each task. High-quality outputs from these specialist models were selected and used to fine-tune the \enquote{base reasoner} to create a \enquote{generalist} model, capable of reasoning about all the tasks. In the last step, the generalist model undergoes another round of \gls{rl} on the chemical tasks to improve the performance. This method shows that structured problem-solving approaches can significantly improve performance on complex chemical tasks without the need for massive domain-specific corpora. 

These diverse approaches illustrate the evolving landscape of chemical \glspl{gpm}.
Still, most applications of \glspl{gpm} focus on using such models for one specific application, and we will review those in the following.

\subsubsection{Limitations}
Chemical \glspl{gpm} are emerging, but we lack systematic understanding of how to build effective ones.

The tokenization question remains unresolved. Domain-specific tokenizers might improve data efficiency, but they prevent reuse of models trained on general text datasets. The tradeoff remains unclear.

A deeper challenge exists: we might lack sufficient high-quality data. Existing chemical datasets are much smaller than those used in domains like mathematics (\Cref{sec:data-section}). The literature reports only successes. Failed experiments, discarded results, and negative outcomes never appear in published datasets. More critically, chemical practitioners rely on tacit knowledge\autocite{Polanyi_2009}---expertise acquired through experience that experts cannot fully articulate. This implicit understanding remains absent from existing datasets.

% Multi-task trained models suggest that performance could benefit from bundling chemistry tasks, yet few studies isolate which objectives or tasks help or harm under different data regimes. Text-based treatment of \gls{smiles} and \gls{cif} overlooks crucial inductive biases, such as symmetries and spatial information, and generations often require post-processing filters to ensure validity and stability. Embeddings from the literature are biased towards the corpus, privileging well-documented chemistries or classes of compounds while ignoring under-explored regions of the chemical space. Scaling the domain pre-training can broaden the coverage, but it does not ensure novelty or robustness in data-starved parts of the chemical space. \gls{rl}-based reasoning adds verification and backtracking, but risks reward-hacking if the validators are weak proxies instead of rigorous simulations or experiments. Evaluations remain benchmark-bound and are fragmented across datasets and metrics; data leakage from pre-training corpora into test sets is seldom audited.

% \subsubsection{Open Challenges}
% \begin{itemize}
% \item \textbf{Data Availability:} Current corpora used for pre-training chemical \glspl{gpm} do not span the full scope and diversity of domains in chemistry and materials science. Moreover, they do not represent the tacit knowledge in these fields.

% \item \textbf{Representations:} A universal structure-aware chemical representation that is amenable to proper tokenization remains elusive. 

% \item \textbf{Chemical Validity and Diversity:} The chemical validity, stability, and synthesizability of the generated candidates are not guaranteed. Furthermore, there needs to be better measures of novelty in composition/structure space than just the distances in embedding space.
% \end{itemize}

\subsection{Knowledge Gathering}\label{sec:information_gathering}
The rate of publishing keeps growing, and as a result, it is increasingly challenging to manually collect all relevant knowledge, potentially stymying scientific progress.\autocite{schilling2025text, Chu_2021}
Even though knowledge collection might seem like a simple task, it often involves multiple steps, visually described in \Cref{fig:knowledge_gathering}A. 
Here, we focus on structured data extraction and question answering. Example queries for both sections are in \Cref{fig:knowledge_gathering}B.


\paragraph{Semantic Search} A step that is key to most, if not all, knowledge-gathering tasks is \gls{rag}, discussed in more detail in \Cref{sec:rag}. 
Most commonly, this involves semantic search, intended to identify chunks of text with similar meaning. 
The difference between semantic search and conventional search lies in how each approach interprets queries. The latter operates through lexical matching---whether exact or fuzzy---focusing on the literal words and their variations. Semantic search, however, focuses on the underlying meaning and contextual relationships within the content. 

To enable semantic search, documents are converted into embeddings (see \Cref{sec:embeddings}) \autocite{bojanowski2017enriching}. 
They allow for similarity search by vector comparison (e.g., using cosine similarity for small databases or more sophisticated algorithms like \gls{hnsw} for large databases\autocite{malkov2018efficient}). 

In chemistry, semantic search has been used extensively to classify and identify chemical text.\autocite{Guo2021,beltagy2019scibert0,trewartha2022quantifying}

\begin{figure}[H]
    \centering
\includegraphics[width=1\textwidth]{figures/final_figures/chemrev_figure12.pdf}
    \caption{\textbf{A. A representation of a typical agent for scientific queries.} The \gls{llm} is the central piece of the system, surrounded by typical tools that improve its question-answering capabilities, together forming an agentic system. The tools represented in this figure are semantic search, citation traversal, evidence gathering, and question answering. Semantic search finds relevant documents. Evidence gathering ranks and filters chunks of text using \glspl{llm}. The citation traversal tool provides model access to citation graphs, enabling accurate referencing of each chunk and facilitating the discovery of additional sources. Finally, the question-answering tool (an \gls{llm}) collects all the information found by other tools and generates a final response to a user's query. This part of the figure is inspired by the \modelname{PaperQA2} agent.\autocite{skarlinski2024language} \textbf{B.} Two examples of applications are discussed in this section. \textbf{C.} worked example using the \texttt{FutureHouse}'s platform (the Crow agent). \enquote{used in 1.1} indicates that this reference has been used in the final response of the model\autocite{skarlinski2025futurehouse}}
\label{fig:knowledge_gathering}
\end{figure}

\subsubsection{Structured Data Extraction}

For many applications, it can be useful to collect data in a structured form, e.g., tables with fixed columns.
Obtaining such a dataset by extracting data from the literature using \glspl{llm} is currently one of the most practical avenues for \glspl{llm} in the chemical sciences \autocite{schilling2025text}.

\paragraph{Data Extraction Using Prompting} 

For most applications, zero-shot prompting should be the starting point. Zero-shot prompting has been used to extract data about organic reactions\autocite{rios2025llm,vangala2024suitability, Patiny2023automatic}, synthesis procedures for metal-organic frameworks\autocite{zheng2023chatgpt}, polymer data\autocite{schilling2024using,gupta2024data}, or other materials data\autocite{polak2024extracting,hira2024reconstructing,kumar2025mechbert,wu2025large,huang2022batterybert}. 
 
\paragraph{Fine-tuning Based Data Extraction} 

If a (commercial) model needs to be run very often, it can be more cost-efficient to fine-tune a smaller, open-source model compared to prompting a large model (see \Cref{sec:distillation}). 
In addition, models might lack specialized knowledge and might not follow certain style guides, which can be introduced with fine-tuning. 
\textcite{ai2024extracting} fine-tuned the \modelname{LLaMa-2-7B} model to extract chemical reaction data from \gls{uspto} into a \gls{json} format compatible with the schema of \gls{ord}\autocite{Kearnes_2021}, achieving an overall accuracy of more than $90\%$. 
In a different approach, \textcite{zhang2024fine} fine-tuned a closed-source model --- \modelname{GPT-3.5-Turbo} to recognize and extract chemical entities from \gls{uspto}. Fine-tuning improved the performance of the base model on the same task by more than $15\%$. 
\textcite{dagdelen2024structured} went beyond the afore-mentioned methods by using a human-in-the-loop data annotation process. Here, humans corrected the outputs from an \gls{llm} extraction instead of annotating data from scratch.

\paragraph{Agents for Data Extraction} 

Agents (\Cref{sec:agents}) have shown their potential in data extraction, though to a limited extent.\autocite{chen2024autonomous,kang2024chatmof}
For example, \textcite{ansari2024agent} introduced \modelname{Eunomia}, an agent that autonomously extracts structured materials data from scientific literature without requiring fine-tuning. Their agent is an \gls{llm} with access to tools such as chemical databases (e.g., the \modelname{Materials Project} database) and research papers from various sources. The document search tool leverages an \gls{api}-based embedding model to find the top-k most relevant passages based on semantic similarity. Other tools used include the chain-of-verification tool, which queries the agent independently (to avoid bias) in order to generate verification queries and finally produces a verified response.

While the authors claim this approach simplifies dataset creation for materials discovery, the evaluation is limited to a narrow set of materials science tasks (mostly focusing on \glspl{mof}), indicating the need for the creation of agent evaluation tools.


\subsubsection{Question Answering}
Besides extracting information from documents in a structured format, \glspl{llm} can also be used to answer questions---such as \enquote{Has X been tried before} by synthesizing knowledge from a corpus of documents (and potentially automatically retrieving additional documents). 

An example of a system that can do that is \modelname{PaperQA}. This agentic system contains tools for search, evidence-gathering, and question answering as well as for traversing citation graphs, which are shown in \Cref{fig:knowledge_gathering}. The evidence-gathering tool collects the most relevant chunks of information via the semantic search and performs \gls{llm}-based re-ranking of these chunks (i.e., the \gls{llm} changes the order of the chunks depending on what is needed to answer the query).
Subsequently, only the top-$n$ most relevant chunks are kept. To further ground the responses, citation traversal tools (e.g., Semantic Scholar\autocite{kinney2023semantic}) are used. 
These leverage the citation graph as a means of discovering supplementary literature references. Ultimately, to address the user's query, a question-answering tool (a specially prompted \gls{llm}) is employed. It initially augments the query with all the collected information before providing a definitive answer.
The knowledge aggregated by these systems could be used to generate new hypotheses or challenge existing ones. 

\subsubsection{Limitations} 

\Glspl{llm} and \gls{llm}-based systems are valuable information-gathering tools, but they have critical limitations. Their knowledge becomes outdated immediately after training. Without additional tools, they cannot identify when retrieved sources have been retracted or corrected. 

Effective knowledge gathering requires human-\gls{ai} collaboration. Current systems struggle to ask relevant clarifying questions.\autocite{choudhury2025bed0llm0} Domain-specific benchmarks for evaluating this capability remain absent, hindering progress in developing truly interactive knowledge-gathering agents.

\subsubsection{Open Challenges}

Critical challenges remain unresolved:

\begin{itemize}
    \item \textbf{Multimodal Understanding.} Current models cannot reliably extract data from plots, tables, and figures.\autocite{alampara2024probing,gupta2022discomat0} This severely limits data extraction because substantial chemical knowledge exists only in visual formats. Improved multimodal capabilities are essential for comprehensive literature mining.
    
    \item \textbf{Source Quality Assessment.} Selecting high-quality sources, that are also relevant to the scientific query, remains an open challenge, with scholarly metrics alone being insufficient.
\end{itemize}


\input{sections/hypothesis}

\subsection{Experiment Planning}
\label{sec:planning}
Before a human or robot can execute any experiments, a plan must be created. 
Planning can be formalized as the process of decomposing a high-level task into a structured sequence of actionable steps aimed at achieving a specific goal. 
The term planning is often confused with scheduling and \gls{rl}, which are closely related but distinct concepts. Scheduling is a more specific process focused on the timing and sequence of tasks. 
It ensures that resources are efficiently allocated, experiments are conducted in an optimal order, and constraints (such as lab availability, time, and equipment) are respected.\autocite{kambhampati2023llmplanning} 
\gls{rl} is about adapting and improving plans over time based on ongoing results.\autocite{chen2022deep}



\subsubsection{Conventional Planning} 

Early chemical planning systems, such as \gls{lhasa} \autocite{corey1972computer} and \modelname{Chematica}\autocite{grzybowski2018chematica}, relied on simple rules and templates. In particular, \modelname{Chematica} used heuristic-guided graph search with rule-based transforms and scoring functions to prune and prioritize routes. Modern systems, like ASKCOS\autocite{tu2025askcos}, explicitly use search algorithms such as \gls{bfs}, \gls{mcts}\autocite{segler2017towards} to explore the combinatorially large space. But these planning search algorithms remain inefficient for long-horizon or complex planning tasks. \autocite{liu2024entropy, zhao2024efficient}


\subsubsection{LLMs to Decompose Problems into Plans}
\glspl{gpm}, in particular \glspl{llm}, can potentially assist in planning with two modes of thinking. 
Deliberate thinking can be used to score potential options or to decompose problems into plans. 
Intuitive thinking can be used to efficiently prune search spaces. 
These two modes align with psychological frameworks known as system-1 (intuitive) and system-2 (deliberate) thinking. \autocite{kahneman2011thinking}
In the system-1 thinking, \glspl{llm} support rapid decision-making by leveraging heuristics and pattern recognition to quickly narrow down options. 
In contrast, system-2 thinking represents a slower, more analytical process, in which \glspl{llm} solve complex tasks by explicitly generating step-by-step reasoning. \autocite{ji2025test}

\Cref{fig:planning_example} shows how a \gls{llm} applies this deliberate, system-2-style reasoning to decompose a chemical problem into a sequence of planned steps.
A variety of strategies have been proposed to improve the reasoning capabilities of \glspl{llm} during inference. 
Methods such as \gls{cot} and least-to-most prompting guide models to decompose. 
However, their effectiveness in planning is limited by error accumulation and linear thinking patterns.\autocite{stechly2024chain}
To address these limitations, recent test-time strategies such as repeat sampling and tree search have been proposed. 
Repeated sampling allows the model to generate multiple candidate reasoning paths, encouraging diversity in thought and increasing the chances of discovering effective subgoal decompositions. \autocite{wang2024planning}
Meanwhile, tree search methods like \gls{tot} and \gls{rap} treat reasoning as a structured search, using algorithms like \gls{mcts} to explore and evaluate multiple solution paths, facilitating more global and strategic decision-making. \autocite{hao2023reasoning}

\glspl{llm} have also been applied to generate structured procedures from limited observations. For example, in quantum physics, a model was trained to infer reusable experimental templates from measurement data, producing Python code that generalized across system sizes. \autocite{arlt2024meta0designing} 

\begin{figure}[!htbp]
    \centering
        \includegraphics[width=1\textwidth]{figures/rescaled_figures/Chemrev_planning_example.pdf}
    \caption{\textbf{An example of using \gls{llm} for chemical planning to decompose problem}. \gls{llm} decomposes a chemical problem into sequential steps with detailed procedures. We manually evaluate the plan step by step, and highlight the problematic steps in text boxes, each linked to the corresponding reason. Plan was generated by ChatGPT-5.
}
    \label{fig:planning_example}
\end{figure}


\subsubsection{Pruning of Search Spaces}

Pruning refers to the process of eliminating unlikely or suboptimal options during the search to reduce the computational burden.
Classical planners employ heuristics, value functions, or logical filters to perform pruning\autocite{bonet2012action}. 

\Cref{fig:planning} illustrates how \glspl{llm} can support experimental planning by pruning options by emulating an expert chemist's intuition by discarding synthetic routes that appear unnecessarily long, inefficient, or mechanistically implausible. 

To further enhance planning efficacy, \glspl{llm} can be augmented with external tools that estimate the feasibility or performance of candidate plans, enabling targeted pruning of the search space before costly execution. 

Beyond external tools, \glspl{llm} can self-correct by pruning flawed reasoning steps to produce more coherent plans. At a higher level of oversight, human-in-the-loop frameworks such as \modelname{ORGANA}\autocite{darvish2025organa} incorporate expert chemist feedback to refine goals, resolve ambiguities, and eliminate invalid directions. Example~\ref{example:chemcrow_organa} presents examples illustrating how planning extends to real laboratory practice.

\begin{figure}[!htbp]
    \centering
        \includegraphics[width=1\textwidth]{figures/rescaled_figures/chemrev_figure14.pdf}
    \caption{\textbf{\gls{gpm}-guided retrosynthesis route planning and pruning}. \glspl{gpm} can systematically evaluate and prune retrosynthetic routes using multiple reasoning capabilities to discriminate between viable and problematic approaches. The partially overlapping arrows at the start of each route indicate multiple steps. \textbf{Route A}: This route was pruned by heuristic reasoning due to the unfavorable aromatic core construction.
    \textbf{Route B}: This route was selected as it successfully passes all \gls{gpm} planning checks, demonstrating optimal synthetic feasibility.
    \textbf{Route C}: This pathway was pruned by external tools due to the poor regio-selectivity of the oxidation step.
    \textbf{Route D}: This route was pruned based on learned intuition, as it represents an inefficient multistep pathway; the route could just start with phenol instead of synthesizing it.
}
    \label{fig:planning}
\end{figure}


\subsubsection{Limitations} 

\Cref{fig:planning_example} demonstrates how \glspl{llm} can generate sophisticated-looking plans that conceal critical flaws, making complex, long-horizon chemical planning tasks particularly difficult to execute reliably \autocite{Cao2025LargeLM}. 
A limitation is that they often produce outputs that appear chemically plausible but are invalid or unsafe, due to their optimization for linguistic plausibility rather than chemical correctness and their lack of mechanistic understanding\autocite{Bran2023TransformersAL, Evans2021TruthfulAD}
Second, errors can also propagate from external tools like retrosynthesis planners, whose data and algorithmic shortcomings constrain reliability \autocite{Li2025ChemHASHA}. 
Finally, a broader limitation is the lack of grounding in experimental feedback, which creates persistent gaps between theoretical planning and practical feasibility.

\subsubsection{Open Challenges}

\begin{itemize}
    \item \textbf{Reasoning Complexity Beyond Knowledge Retrieval}: Complex chemistry problems require long, tightly interconnected chains of reasoning, where minor errors can cascade and require understanding of dynamic interactions such as temperature effects on molecular behavior. Current \glspl{llm} lack effective reasoning structures to guide domain-specific reasoning. \autocite{Ouyang2023StructuredCR, Tang2025ChemAgentSL} 
    \item \textbf{Evaluation and Feedback Bottleneck}: Current evaluation methods are often performed manually or indirectly, either relying on expert review as in \modelname{ChemCrow} \autocite{bran2024augmenting} or on pseudocode-based comparisons as in \modelname{BioPlanner}.\autocite{o2023bioplanner} Integrating feedback remains an open direction for improving the practical feasibility of generated plans.
\end{itemize}


\subsection{Experiment Execution}
Once an experimental plan is available, whether from a human scientist's idea or an AI model, the next step is to execute it. 
Regardless of its source, the plan must be translated into concrete, low-level actions for execution. 

It is worth noting that, despite their methodological differences, executing experiments \textit{in silico} (running simulations or code) and \textit{in vitro} are not fundamentally different---both follow an essentially identical workflow: Plan $\rightarrow$ Instructions $\rightarrow$ Execution $\rightarrow$ Analysis. 
% In a computer simulation, a researcher writes a program (plan), which is then compiled or interpreted into machine code (instructions) for the \gls{cpu}, executed to produce data, and finally, the outputs are analyzed. 
% In an automated laboratory, the scientist specifies a protocol (plan), which must be translated into instrument commands (instructions), executed on a robotic platform, followed by the analysis of sensor data or assay results.

The execution of \textit{in silico} experiments can be reduced to two essential steps: preparing input files and running the computational code; \glspl{gpm} can be used in both steps.\autocite{Liu2025ASA, Mendible‑Barreto2025DynaMate, Zou2025ElAgente, Campbell2025MDCrow} \textcite{Jacobs2025orca} found that using a combination of fine-tuning, \gls{cot} and \gls{rag} (see \Cref{sec:model_adaptation}) can improve the performance of \glspl{llm} in generating executable input files for the quantum chemistry software \emph{ORCA}\autocite{ORCA5}, while \textcite{Gadde2025chatbot} created \modelname{AutosolvateWeb}, an \gls{llm}-based platform that assists users in preparing input files for \gls{qmmm} simulations of explicitly solvated molecules and running them on a remote computer. 
Examples of \gls{llm}-based autonomous agents (see 
\Cref{sec:agents}) capable of performing the entire computational workflow (i.e., preparing inputs, executing the code, and analyzing the results) are \modelname{MDCrow} \autocite{Campbell2025MDCrow} (for molecular dynamics) and \modelname{El Agente Q} \autocite{Zou2025ElAgente} (for quantum chemistry).

Emerging examples show \glspl{gpm} assisting in \textit{in vitro} experiment automation. Programming language paradigms---compiled vs.\ interpreted (\Cref{fig:exec}A)---provide a useful analogy for understanding different automation approaches.

Compiled languages (\modelname{C++}, \modelname{Fortran}) convert entire programs to machine code before execution. Interpreted languages (\modelname{Python}, \modelname{JavaScript}) translate and execute instructions line-by-line at runtime. The tradeoff is that compiled languages offer higher performance and early error detection but require separate compilation steps. Interpreted languages enable rapid development and on-the-fly modification, but run slower and catch errors only during execution.

Similarly, experiment automation follows two paradigms (\Cref{fig:exec}B): \enquote{compiled automation} translates entire protocols---by human or \gls{gpm}---into low-level instructions before execution. \enquote{Interpreted automation} uses the \gls{gpm} as runtime interpreter, executing protocols step-by-step.


\begin{figure}[H]
    \centering
\includegraphics[width=\textwidth]{figures/rescaled_figures/chemrev_figure15+16.pdf}
    \caption{\textbf{Programming languages vs.\ lab automation. A) programming paradigms}: In compiled languages, the entire source code is translated ahead of time to machine code by the compiler. This stand-alone code is then given to the \gls{os}, which is responsible for scheduling and distributing tasks to the hardware. In interpreted languages, the interpreter reads and translates each line of the source code to machine code and hands it to the \gls{os} for execution. \textbf{B) automation paradigms}: In the compiled approach, a \gls{gpm} formalizes the protocol, a compiler, such as the Chempiler\autocite{steiner2019organic}, translates the formalized protocol to hardware-specific low-level steps, which the controller then executes---a central hub tasked with scheduling and distributing commands to chemical hardware. In the interpreted approach, a \gls{gpm}, acting as the interpreter, first breaks down the protocol into specific steps, then sends them (via an \gls{api}) for execution one by one. The strength of interpreted systems is dynamic feedback: after the execution of each step, the \gls{gpm} receives a signal (e.g., data, errors), which can influence its behavior for the next steps.}
    \label{fig:exec}
\end{figure}

\subsubsection{Compiled Automation}
In \enquote{compiled automation}, protocols are formalized in high-level languages or \glspl{dsl}. A chemical compiler (\enquote{chempiler}\autocite{mehr2020universal}) converts these into low-level hardware instructions, which robotic instruments then execute (\Cref{fig:exec}B). 

While \modelname{Python} scripts serve as the \textit{de facto} protocol language, specialized \glspl{dsl} provide more structured representations.\autocite{wang2022ulsa, ananthanarayanan2010biocoder, autoprotocol2023, Park2023CMDL} For example, \gls{chidl}\autocite{mehr2020universal, xdl2023spec} describes protocols using abstract commands (\modelname{Add}, \modelname{Stir}, \modelname{Filter}) and chemical objects (\modelname{Reagents}, \modelname{Vessels}). The \modelname{Chempiler} translates \gls{chidl} scripts into platform-specific instructions based on the physical laboratory configuration.

Writing protocols in formal languages requires coding expertise. Here, \glspl{gpm} translate natural-language protocols into machine-readable formats.\autocite{Lamas2024DSLXpert, jiang2024protocode, conrad2025lowering, inagaki2023robotic,Vaucher2020AutoExtraction}
\textcite{Pagel2024LLMChemputer} introduced a multi-agent workflow (based on \modelname{GPT-4}) that can convert unstructured chemistry papers into executable code. 
The first agent extracts all synthesis-relevant text, including supporting information; a procedure agent then sanitizes the data and tries to fill the gaps from chemical databases (using \gls{rag}); another agent translates procedures into \gls{chidl} and simulates them on virtual hardware; finally, a critique agent cross-checks the translation and fixes errors.

The example above shows one of the strengths of the compiled approach: it allows for pre-validation. 
The protocol can be simulated or checked for any errors before running on the actual hardware, ensuring safety. 
Another example of \gls{llm}-based validators for chemistry protocols is \modelname{CLAIRify}.\autocite{Yoshikawa2023CLAIRify} 
Leveraging an iterative prompting strategy, it uses \modelname{GPT‑3.5} to first translate the natural-language protocol into \gls{chidl} script, then automatically verifies its syntax and structure, identifies any errors, appends those errors to the prompt, and prompts the \gls{llm} again---iterating this process until a valid \gls{chidl} script is produced. 


\subsubsection{Interpreted Automation}
Interpreted programming languages support higher abstraction levels through flexible command structures. Similarly, \glspl{gpm} can translate high-level goals into concrete steps,\autocite{ahn2022can, huang2022language} acting as \enquote{interpreters} between experimental intent and lab hardware.

For example, given \enquote{titrate the solution until it turns purple}, a \gls{gpm} agent (\Cref{sec:agents}) can break this into executable steps at runtime: add titrant incrementally, read color sensor, loop until condition met. This is \enquote{interpreted automation}---conversion happens during execution, not before.

The key advantage is real-time decision-making. After each action, the system analyzes sensor data (readings, spectra, errors) and selects the next step. This enables dynamic branching and conditional logic impossible in pre-compiled protocols.

\modelname{Coscientist}\autocite{boiko2023autonomous} demonstrates interpreted automation using \modelname{GPT-4} to control liquid-handling robots. The system searches the web for protocols, reads instrument documentation, writes \modelname{Python} code in real-time, and executes experiments on physical hardware. When errors occur, \modelname{GPT-4} debugs its own code. \modelname{Coscientist} successfully optimized palladium cross-coupling reactions, outperforming Bayesian optimization in finding high-yield conditions.

\modelname{ChemCrow}\autocite{bran2024augmenting} augments \modelname{GPT-4} with $18$ expert-designed tools for tasks including compound lookup, spectral analysis, and retrosynthesis. It planned and executed syntheses of \gls{deet} and three thiourea organocatalysts (example~\ref{example:chemcrow_organa}), and collaborated with chemists to discover a new chromophore.


\subsubsection{Hybrid Approaches}
Between fully compiled and fully interpreted automation lies a hybrid approach that combines the safety and reliability of compiled protocols with the AI-driven flexibility of interpreted systems. Each experiment run follows a fixed plan for safety and reproducibility, but between runs, the plan can adapt based on the GPM’s interpretation of results. This design provides a safeguard against interpreter errors, since every generated procedure passes through formal verification before execution---catching issues like a hallucinated instruction to add \SI{1000}{\milli\liter} of solvent to a \SI{100}{\milli\liter} flask.


\modelname{ORGANA} \autocite{darvish2025organa} is an \gls{llm}-based robotic assistant following this hybrid paradigm. 
It allows human chemists to describe their experimental goal in natural language. The system can converse with the user to clarify ambiguous requests (the agent would ask \enquote{do you mean X or Y?} if the instructions are unclear). Once the goal is understood, it uses \modelname{CLAIRify} \autocite{Yoshikawa2023CLAIRify} to convert and validate the natural-language description of a chemistry experiment into a \gls{chidl} script, which can be executed on a compatible platform. 
In one case, \modelname{ORGANA} carried out a multistep electrochemistry procedure---polishing electrodes, running an experiment, and analyzing the data---involving 19 substeps that it coordinated in parallel (see example~\ref{example:chemcrow_organa}).
If an unexpected observation occurred (e.g., a solution does not change color when expected), the system can notice via image analysis and modify the plan or alert the user. 
In user studies, \modelname{ORGANA} reduced the manual labor and frustration for chemists, who could offload tedious tasks and trust the agent to handle low-level decisions.


\begin{examplebox}[label={example:chemcrow_organa}]{Planning and executing experiments with \glspl{gpm}}

\textbf{ChemCrow:} When prompted with the task \enquote{Find and synthesize a thiourea organocatalyst which accelerates a Diels-Alder reaction}, \modelname{ChemCrow} sequentially queried tools to gather information and solve the task:

\par\vspace{0.5\baselineskip}

{\leftskip=2em \enquote{First, I need to find a thiourea organocatalyst that can accelerate the Diels-Alder reaction. I will perform a websearch to find a suitable catalyst.

\textbf{WebSearch tool:} Schreiner's thiourea catalyst

Now I will obtain the \gls{smiles}. 

\textbf{Name2Smiles tool:}

\texttt{FC(F)(F)c1cc(CN(=S)Nc2cc(C(F)(F)F)cc(C(F)(F)F)c2)cc(C(F)(F)F)c1}

I will plan a synthesis for Schreiner's thiourea catalyst. 

\textbf{SynthesisPlanner tool:} Detailed synthesis plan

\includegraphics[width=0.8\textwidth]{figures/schreiner.png}

I will execute the synthesis.

\textbf{SynthesisExecuter tool:} Successful
synthesis.}

}

\par\vspace{0.5\baselineskip}

The agent proposed to synthesize the target molecule from the corresponding isothiocyanate and aniline derivatives and executed the synthesis on the RoboRXN platform. \autocite{roborxn} The product was confirmed using HPLC/MS: \texttt{m/z 501 [M+H]} calculated, found \texttt{m/z 501.02}.

\par\vspace{\baselineskip}

\textbf{ORGANA}: \modelname{ORGANA} interacts with chemists using \glspl{llm} to derive experiment goals and makes decisions with chemists in the loop to control robots and lab devices.
\begin{itemize}
    \item \textbf{Start-up:} Users give their intention \enquote{I want to create a Pourbaix plot of a quinone, Sodium Anthraquinone-1-sulfonate.}, and start interact with \modelname{ORGANA}:
    
    \modelname{ORGANA}: What is the setup of the lab environment? What hardware and reagents are present?
    
    User: I have a \SI{2}{\milli\molar} solution of the quinone......
    
    \modelname{ORGANA}: Can you provide an example of how to run the experiment? 
    
    User: Add \SI{2}{\milli\liter} of quinone......

    \modelname{ORGANA}: What is the expected outcome of the example experiment? 
    
    User: The pH of the solution is expected ......
    
    \modelname{ORGANA}: How many experiments are we planning to run in total? 
    
    User: 6
    
    \modelname{ORGANA}: Ok, thank you! I will start the experiment......

    \item \textbf{Ambiguity and Uncertainty Mitigation:} 
    \modelname{ORGANA.Reasoner} evaluates whether experiment outcomes align with user expectations. If it can’t interpret the results, it alerts the user to check for environmental issues. The user provides feedback explaining the outcomes, which \modelname{ORGANA.Reasoner} then uses to refine its next plan:
    
    \enquote{I’m not sure what happened, but repeat the previous experiment just in case}
    
    \enquote{Nothing is wrong, carry on}

    \enquote{The pump was stuck, it’s ok again}
    
    
    \item \textbf{Experimental Reports:} 
    \modelname{ORGANA} reports the experimental logs and summaries:
    
    \enquote{A series of experiments was conducted to measure the potential of a quinone solution at various pH levels, specifically from $pH \space 7$ to $pH \space 9$......
    After performing the experiments, these are the results:
    The estimate for $pK_{a1}$ is $8.096$.
    The estimate for $pK_{a2}$ is $12.380$.
    The estimate for slope is $-60.958$.}

    \item \textbf{Results Comparison with Chemists:} Comparison of Pourbaix diagrams and first-region slope estimates from electrochemical experiments, with \modelname{ORGANA} using three measurement points per pH region and chemists using four. \modelname{ORGANA} yields results comparable to those of chemists: for $pK_{a1}$, \modelname{ORGANA} produces $8.03$ and chemists $8.02$ and for the estimated slope, \modelname{ORGANA} obtained \SI[per-mode=symbol]{-61.3}{\milli\volt\per\pH} while chemists obtained \SI[per-mode=symbol]{-62.7}{\milli\volt\per\pH}.

\end{itemize}

\end{examplebox}

\subsubsection{Limitations}
Current automation systems remain prototypes. 

Interpreted systems require frequent human intervention despite autonomy claims. They replicate known procedures but lack mechanistic understanding. Non-deterministic \gls{gpm} responses create reproducibility issues---small prompt changes yield different results, and closed-source models evolve unpredictably. Hallucinations risk incorrect planning for complex reactions. Hardware control introduces safety concerns: flexible \glspl{gpm} can devise unanticipated actions, requiring robust safeguards (\Cref{sec:safety}).

Compiled systems offer reliability but require extensive upfront formalization. The effort to translate protocols into formal languages often outweighs automation benefits for typical laboratory workflows.

%\Cref{tab:execution_comparison} compares the strengths and weaknesses of the different experiment execution paradigms discussed. 

% \begin{table}[ht]
% \centering
% \caption{\textbf{Comparison of the different paradigms of lab automation used alongside \glspl{gpm}}. Each approach has its strengths and weaknesses. Compiled systems favor reliability, interpreted systems allow for more flexibility, while hybrid systems try to strike a balance. }
% \begin{tabular}{lccc}
% \toprule
% \textbf{Feature} & \textbf{Compiled} & \textbf{Interpreted} & \textbf{Hybrid} \\
% \midrule
% Flexibility & \textcolor{NegativeColor}{Low} & \textcolor{PositiveColor}{High} & Medium \\
% Adaptivity & \textcolor{NegativeColor}{None} & \textcolor{PositiveColor}{Real-time} & Iterative \\
% Reproducibility & \textcolor{PositiveColor}{High} & Medium & \textcolor{PositiveColor}{High} \\
% Safety & \textcolor{PositiveColor}{High} & \textcolor{NegativeColor}{Low} & Medium \\
% Setup Overhead & Medium & \textcolor{NegativeColor}{High} & \textcolor{NegativeColor}{High} \\
% Industrial Readiness & \textcolor{NegativeColor}{Low} & \textcolor{NegativeColor}{Low} & \textcolor{NegativeColor}{Low} \\
% \bottomrule
% \end{tabular}
% \label{tab:execution_comparison}
% \end{table}

\subsubsection{Open Challenges}
Self-driving laboratories orchestrated by \glspl{gpm} face technical challenges requiring research advances:\autocite{Tom2024SDL,Seifrid2022SDL}

\begin{itemize}
\item \textbf{Grounding Natural Language to Laboratory Actions.} Translating ambiguous natural-language instructions (\enquote{heat gently}) into precise, safe operations requires developing validation layers that detect physically impossible or hazardous actions before hardware execution.

\item \textbf{Universal Protocol Standards.} No widely adopted formalization standard exists. While languages like \gls{chidl} show promise, achieving interoperability across platforms requires community consensus on abstraction levels and device interfaces. \Glspl{mcp} offer a potential path forward by enforcing consistent interfaces between \glspl{gpm}, instruments, and verification layers.

\item \textbf{Autonomous Error Recovery.} Current systems cannot autonomously diagnose and recover from experimental failures. Developing general-purpose failure detection mechanisms and recovery strategies would enable truly autonomous operation.

\item \textbf{Multimodal Integration.} Chemists use diverse data types---spectra, chromatograms, TLC plates, and microscopy images. Integrating these modalities into \gls{gpm} decision-making loops remains technically challenging but essential for human-level experimental reasoning.

\item \textbf{Verification and Provenance.} Industrial and clinical applications require complete experimental provenance: every decision logged with reasoning traces, all parameters recorded, and outcomes traceable to specific model versions (\Cref{sec:safety}).
\end{itemize}

\subsection{Data Analysis}
The analysis of experimental data in chemistry remains a predominantly manual process. 

One key challenge that makes automation particularly difficult is the extreme heterogeneity of chemical data sources. 
Laboratories often rely on a wide variety of instruments, some of which are decades old, rarely standardized, or unique in configuration.\autocite{jablonka2022making} 
These devices output data in incompatible, non-standardized, or poorly documented formats, each requiring specialized processing pipelines. 
Despite efforts like \modelname{JCAMP-DX} \autocite{McDonald1988standard}, standardization attempts remain limited and have generally failed to gain widespread use. 
This diversity makes rule-based or hard-coded solutions largely infeasible, as they cannot generalize across the long tail of edge cases and exceptions found in real-world workflows.

This complexity makes chemical data analysis a promising application for \glspl{gpm} (\Cref{fig:anaylsis}). These models handle diverse tasks and formats using implicit knowledge from broad training data. In other domains, \glspl{llm} have successfully processed heterogeneous tabular data and performed classical data analysis without task-specific training.\autocite{narayan2022can, kayali2023chorus}

\begin{figure}[!ht]
    \centering
\includegraphics[width=1\textwidth]{figures/rescaled_figures/chemrev_figure17.pdf}
    \caption{\textbf{Static conventional data analysis workflow vs.\ dynamic \gls{gpm} generated workflow}. Chemical analysis can be performed with a variety of instruments and techniques, resulting in many possible output data formats. The \gls{gpm} can use these diverse, raw data and process them into easy-to-understand plots, analysis, and reports. A hard-coded workflow, in contrast, is specifically made to analyze one specific data format and spectra and produces a fixed output format, e.g., the \gls{smiles} of the analyzed molecule.}
    \label{fig:anaylsis}
\end{figure}

In chemistry, however, only a handful of studies have so far demonstrated similar capabilities. Early evaluations showed that \glspl{gpm} can support basic workflows such as the classification of \gls{xps} signals based on peak positions and intensities.\autocite{Fu2025large,decurt2024large} 
 
Spectroscopic data often appear as raw plots or images, making direct interpretation by \glspl{vlm} a more natural starting point for automated analysis. 
A broad assessment of \gls{vlm}-based spectral analysis was introduced with the \modelname{MaCBench} benchmark \autocite{alampara2024probing}, which systematically evaluates how \glspl{vlm} interpret experimental data in chemistry and materials science---including various types of spectra such as \gls{ir}, \gls{nmr}, and \gls{xrd}---directly from images. They showed that while \glspl{vlm} can correctly extract isolated features from plots, the performance substantially drops in tasks requiring deeper spatial reasoning.
To overcome these limitations, \textcite{kawchak2024high} explored two-step pipelines that decouple visual perception from chemical reasoning. First, the model interprets each spectrum individually (e.g., converting  \gls{ir}, \gls{nmr}, or \gls{ms} images into textual peak descriptions), and second, a \gls{llm} analyzes these outputs to propose a molecular structure based on the molecular formula. 


More complex agentic systems extend beyond single-step analysis and attempt to orchestrate entire workflows. 
For example, \textcite{ghareeb2025robin0} developed a multi-agent system for assisting biological research with hypothesis generation (see \Cref{fig:hypothesis-generation}) and experimental analysis. Its data analysis agent \modelname{Finch} autonomously processes raw or preprocessed biological data, such as \gls{rna} sequencing or flow cytometry, by executing code in Jupyter notebooks and producing interpretable outputs. Currently, only these two data types are supported, and expert-designed prompts are still required to ensure reliable results. 

Similarly, \textcite{mandal2024autonomous} introduced \modelname{AILA}, which utilizes \gls{llm}-agents to plan, execute, and iteratively refine full \gls{afm} analysis pipelines. Compared to earlier prototypes, this systems emphasize transparency and reproducibility by producing both code and reports.
The system handles tasks such as image processing, defect detection, clustering, and the extraction of physical parameters. 

\subsubsection{Limitations}
While \glspl{gpm} offer promising capabilities for automating scientific data analysis, several concrete limitations remain. 
Recent evaluations such as \modelname{SciCode} \autocite{tian2024scicode} have shown that even \gls{sota} like \modelname{GPT-4-Turbo} frequently produce syntactically correct but semantically incorrect code, for instance, in common data analysis steps such as reading files, applying filters, or generating plots.

These technical shortcomings are further amplified by sensitivity to prompt formulation. As demonstrated by \textcite{Yan2020auto} and \textcite{alampara2024probing}, even minor changes in wording or structure can lead to drastically different results, highlighting a lack of robustness in prompt-based control. 

In practice, this means that robust prompting strategies, systematic validation, and human oversight remain essential components of any current deployment.

\subsubsection{Open Challenges}
Looking forward, several open challenges remain unresolved.

\begin{itemize}
    \item \textbf{True Chemical Reasoning}: It remains unclear whether current \glspl{llm} can perform genuine chemical analysis rather than relying on pattern-matching or shallow feature extraction. \autocite{alampara2024probing, alampara2025task}
    \item \textbf{Seamless Laboratory Integration}: No commercial systems yet provide robust, end-to-end interoperability with the diverse analytical instruments used in chemistry laboratories. Existing research prototypes support only limited data types and still depend heavily on expert curation. \autocite{ghareeb2025robin0, mandal2024autonomous}
    \item \textbf{Standardization and Real-World Validation}: Achieving production-ready systems requires progress not only in model robustness but also in data and protocol standardization, hardware integration, and thorough testing under realistic laboratory conditions. \autocite{testini2025measuring}
\end{itemize}


\subsection{Reporting}
To share insights obtained from data analysis, one often converts them into scientific publications or other forms of content, such as reports or blogs.
In this step, \glspl{gpm} can also take a central role. While writing assistance has been showcased in past works, it remains limited in scope and real-world impact.


\paragraph{From Data to Explanation} The lack of explainability of \gls{ml} predictions generates skepticism among experimental chemists\autocite{wellawatte2025human}, hindering the wider adoption of such models.\autocite{wellawatte2022model}
One promising approach to address this challenge is to convey explanations of model predictions in natural language. 
An approach proposed by \textcite{wellawatte2025human} involves coupling \glspl{llm} with feature importance analysis tools, such as \gls{shap} or \gls{lime}. 
In this framework, the \gls{llm} performs three key functions: First, it translates technical feature names into more accessible language. Second, using \gls{rag} over scientific literature, it retrieves relevant excerpts that explain the physicochemical relationships between identified molecular features and target properties. Third, it synthesizes these components into coherent natural language explanations that not only identify which structural features correlate with the property of interest, but also hypothesize \textit{why} these relationships exist based on established chemical principles from the literature. 

\paragraph{Writing Assistance} \glspl{llm} can assist with syntax improvement, redundancy identification,\autocite{khalifa2024using} figure and table captioning,\autocite{hsu2021scicap,selivanov2023medical} caption-figure matching,\autocite{hsu2023gpt04} and alt-text generation.\autocite{singh2024figura11y} Models can be personalized for specific audiences or writing styles.\autocite{li2023teach}

\glspl{llm} has also been shown to potentially help complete submission checklists. \textcite{goldberg2024usefulness} found $70\%$ of \gls{neurips} 2025 authors found \gls{llm} assistance useful for checklist completion, with the same fraction revising their submissions based on model feedback.

\subsubsection{Limitations}
\gls{llm} explanations appear credible but often lack faithfulness to underlying reasoning.\autocite{agarwal2024faithfulness} Models can reinforce existing biases through training data or prompting strategies.\autocite{kobak2025delving} While \glspl{llm} can process large datasets, they miss subtle artifacts and anomalies human researchers detect, and struggle distinguishing correlation from causation.\autocite{jin2023large}

\subsubsection{Open Challenges}
\begin{itemize}
    \item \textbf{Provenance Tracking Systems.} Developing methods to trace every claim back to specific training examples or retrieved sources. This requires architectures that maintain explicit links between generated text and source materials, enabling verification of attribution completeness.
    
    \item \textbf{Authorship Frameworks.} Defining contribution taxonomies that specify when \gls{llm} use constitutes co-authorship versus tool use. Journals and institutions need consensus guidelines for disclosure, attribution, and accountability when \glspl{llm} assist in research reporting.
\end{itemize}